\section*{Подробные решения. Вариант 5}

\subsection*{1. Восьмеричное число в шестнадцатеричной системе}
Заменяем восьмеричные цифры тройками битов, затем группируем по четыре в обе стороны от точки:
\[
\begin{aligned}
7652.34_8&=\texttt{111 110 101 010.011 100}_2\\
&=\texttt{1111 1010 1010.0111 0000}_2\\
&=\boxed{FAA.7_{16}}.
\end{aligned}
\]
Дробная часть: $0.34_8=3/8+4/64=7/16=0.7_{16}$. Дописанные справа нули значение не меняют.

\subsection*{2. Перевод десятичного целого}
Последовательно делим на 16:
\[
\begin{aligned}
9110&=16\cdot569+6, &569&=16\cdot35+9,\\
35&=16\cdot2+3, &2&=16\cdot0+2.
\end{aligned}
\]
Остатки снизу вверх дают $\boxed{2396_{16}}$. Проверка: $2\cdot4096+3\cdot256+9\cdot16+6=9110$.

\subsection*{3. Перевод десятичной дроби}
Каждый шаг умножает оставшуюся дробь на 16; целые части дают цифры:
\[
\begin{array}{rcl@{\qquad}rcl}
0.4486\cdot16&=&7.1776&0.1776\cdot16&=&2.8416,\\
0.8416\cdot16&=&13.4656&0.4656\cdot16&=&7.4496,\\
0.4496\cdot16&=&7.1936&0.1936\cdot16&=&3.0976.
\end{array}
\]
$13=D_{16}$, шестой знак равен 3, поэтому округление до пяти знаков не меняет пятый. Ответ $\boxed{0.72D77_{16}}$.

\subsection*{4. Сложение в системе с основанием 11}
Цифра $A$ означает десять. Дробная часть второго слагаемого равна нулю. Для целых разрядов справа налево получаем $3+9=12=11+1$, $2+9+1=12$, $1+9+1=11$, $A+1=11$. Каждый раз пишем остаток и переносим единицу:
\[
\begin{array}{r@{}l}
A123&.115_{11}\\
\phantom{A}999&.000_{11}\\\hline
10011&.115_{11}
\end{array}
\]
Ответ $\boxed{10011.115_{11}}$.

\subsection*{5. Сложение в 16-битном дополнительном коде}
$342_{10}=0156_{16}$. Для отрицательного $B=-430$: положительное число $430=01AE_{16}$; инверсия 16 битов даёт $FE51_{16}$, прибавление 1 --- $FE52_{16}$. Складываем коды:
\[
0156_{16}+FE52_{16}=\boxed{FFA8_{16}}.
\]
Проверка знакового значения: $342-430=-88$; его код равен $2^{16}-88=65448=FFA8_{16}$. Знакового переполнения нет.

\subsection*{6. Вычитание в дополнительном коде}
$1005_8=512+5=517$, следовательно, $B=-517$ и $A-B=21-(-517)=538_{10}$. Код $A=0015_{16}$; для вычитания отрицательного $B$ прибавляем код $-B=517=0205_{16}$:
\[
0015_{16}+0205_{16}=\boxed{021A_{16}}.
\]
Проверка: $021A_{16}=2\cdot256+16+10=538$. Код самого $B$ --- $FDFB_{16}$.

\subsection*{7. Представление в \texttt{float32}}
$1234.617_{10}$ положительно, значит $S=0$. Из $2^{10}\le1234.617<2^{11}$ следует истинный порядок $E=10$, хранимый $e=10+127=137=10001001_2$. Шаг между представимыми числами равен $2^{10-23}=2^{-13}$. Масштабируем:
\[
1234.617\cdot2^{13}=10113982+\frac{58}{125}.
\]
Так как $58/125<1/2$, округляем до $10113982$. Его 24-битная запись --- \texttt{100110100101001110111110}.
\par\noindent Убираем скрытую старшую единицу. Дробное поле:
\begin{center}
\texttt{00110100101001110111110}.
\end{center}
Соединяем поля:
\[
\texttt{0 10001001 00110100101001110111110}
=\boxed{449A53BE_{16}}.
\]
Код представляет ближайшее к исходному число в \texttt{float32}:
\[
1234.616943359375.
\]

\subsection*{8. Функция логической схемы}
ИЛИ получает $A$ и $B$, исключающее ИЛИ --- $A$ и $C$, нижний НЕ --- $B$. Объединяя результаты И и инвертируя выход, получаем
\[
\boxed{F=\lnot\bigl((A\lor B)\land(A\oplus C)\land\lnot B\bigr)}.
\]
